\setcounter{section}{25}

  \zwischenueberschrift{Übungsaufgaben}

\inputaufgabegibtloesung
{}
{

Es sei $X$ eine
\definitionsverweis {differenzierbare Mannigfaltigkeit}{}{}
und
\maabb {p} { E  =  X \times \R } { X
} {}
das triviale Vektorbündel über $X$ vom
\definitionsverweis {Rang}{}{}
$1$. Es sei $\omega$ eine
$1$-\definitionsverweis {Differentialform}{}{}
auf $X$. Zeige die folgenden Aussagen.
\aufzaehlungfuenf{Die Abbildung
\maabbeledisp {} {TE} { p^*E
} {(P,u,v,w) } { (P, u, \omega(P,v) +w)
} {,}
definiert einen
\definitionsverweis {Zusammenhang}{}{.}
}{Die vertikale Ableitung zu dem Zusammenhang ist für stetig differenzierbare Funktionen
\maabb {f} { U } {\R
} {}
auf einer offenen Menge

\mavergleichskette
{\vergleichskette
{ U
}
{ \subseteq }{ X
}
{  }{
}
{    }{
}
{   }{
}
}
{}{}{}
durch

\mavergleichskettedisp
{\vergleichskette
{ \nabla (f) (P)
}
{ =} { \omega (P) + (df)_P
}
{ } {
}
{ } {
}
{ } {
}
}
{}{}{}
gegeben
\zusatzklammer {hierbei identifizieren wir auf beiden Seiten eine reellwertige Funktion $f$  mit dem Schnitt
\maabbele {} { X } { X \times \R
} { x } { (x,f(x))
} {.}} {} {}
}{Eine stetig differenzierbare Funktion
\maabb {f} { U } {\R
} {}
ist genau dann ein
\definitionsverweis {horizontaler Schnitt}{}{}
über $U$ bezüglich $\nabla$, wenn  $-f$ eine
\definitionsverweis {Stammform}{}{}
zu $\omega$ über $U$ ist.
}{Die Form $\omega$ ist genau dann
\definitionsverweis {geschlossen}{}{,}
wenn der Zusammenhang
\definitionsverweis {lokal integrabel}{}{}
ist.
}{Die Form $\omega$ ist genau dann
\definitionsverweis {exakt}{}{,}
wenn der Zusammenhang
\definitionsverweis {global integrabel}{}{}
ist.
}

}
{} {}

\inputaufgabe
{}
{

Wir betrachten das triviale Vektorbündel
\maabbdisp {p} {E= \R^2 \times \R} { \R^2
} {.}
Beschreibe einen
\definitionsverweis {Zusammenhang}{}{}
auf $E$ derart, dass es keinen
\definitionsverweis {horizontalen Schnitt}{}{}
gibt.

}
{} {}

\inputaufgabe
{}
{

Es sei $X$ eine
\definitionsverweis {differenzierbare Mannigfaltigkeit}{}{}
und
\maabb {p} { E  =  X \times \R } { X
} {}
das triviale Vektorbündel über $X$ vom
\definitionsverweis {Rang}{}{}
$1$. Es sei $\omega$ eine
$1$-\definitionsverweis {Differentialform}{}{}
auf $X$. Zeige die folgenden Aussagen.
\aufzaehlungdrei{Die Abbildung
\maabbeledisp {} {TE} { p^*E
} {(P,u,v,w) } { (P, u, u\omega(P,v) +w)
} {,}
definiert einen
\definitionsverweis {Zusammenhang}{}{.}
}{Die vertikale Ableitung zu dem Zusammenhang ist für stetig differenzierbare Funktionen
\maabb {f} { U } {\R
} {}
auf einer offenen Menge

\mavergleichskette
{\vergleichskette
{ U
}
{ \subseteq }{ X
}
{  }{
}
{    }{
}
{   }{
}
}
{}{}{}
durch

\mavergleichskettedisp
{\vergleichskette
{ \nabla (f) (P)
}
{ =} { f\omega (P) + (df)_P
}
{ } {
}
{ } {
}
{ } {
}
}
{}{}{}
gegeben
\zusatzklammer {hierbei identifizieren wir auf beiden Seiten eine reellwertige Funktion $f$  mit dem Schnitt
\maabbele {} { X } { X \times \R
} { x } { (x,f(x))
} {.}} {} {}
}{Eine nullstellenfreie stetig differenzierbare Funktion
\maabb {f} { U } {\R
} {}
ist genau dann ein
\definitionsverweis {horizontaler Schnitt}{}{}
über $U$ bezüglich $\nabla$, wenn  $-   \ln  \betrag {  f  }$ eine
\definitionsverweis {Stammform}{}{}
zu $\omega$ über $U$ ist.
}

}
{} {}

\inputaufgabe
{}
{

Es sei $X$ eine
\definitionsverweis {differenzierbare Mannigfaltigkeit}{}{}
und $E$ ein
\definitionsverweis {differenzierbares Vektorbündel}{}{}
auf $X$, das mit einem
\definitionsverweis {Zusammenhang}{}{}
versehen sei. Zeige, dass ein stetig differenzierbarer Schnitt
\maabb {s} { U } { E\vert_U
} {}
auf einer offenen Menge

\mavergleichskette
{\vergleichskette
{ U
}
{ \subseteq }{ X
}
{  }{
}
{    }{
}
{   }{
}
}
{}{}{}
genau dann
\definitionsverweis {horizontal}{}{}
ist, wenn seine
\definitionsverweis {vertikale Ableitung}{}{}

\mavergleichskette
{\vergleichskette
{ \nabla s
}
{ = }{ 0
}
{  }{
}
{    }{
}
{   }{
}
}
{}{}{}
ist.

}
{} {}

\inputaufgabe
{}
{

Es sei $X$ eine
\definitionsverweis {differenzierbare Mannigfaltigkeit}{}{}
und $E$ ein
\definitionsverweis {differenzierbares Vektorbündel}{}{}
auf $X$, das mit einem
\definitionsverweis {Zusammenhang}{}{}
versehen sei. Es sei $Z$ eine weitere differenzierbare Mannigfaltigkeit mit einer differenzierbaren Abbildung
\maabb {\psi} { Z } { X
} {.}
Zeige, dass ein stetig differenzierbarer Schnitt
\maabb {s} { Z } { E
} {}
über $\psi$ genau dann
\definitionsverweis {horizontal}{}{}
ist, wenn seine
\definitionsverweis {vertikale Ableitung}{}{}

\mavergleichskette
{\vergleichskette
{ \nabla s
}
{ = }{ 0
}
{  }{
}
{    }{
}
{   }{
}
}
{}{}{}
ist.

}
{} {}

Für die folgende Aufgabe verwende man
[[Mannigfaltigkeit/Vektorbündel/R/Zusammenhang/Horizontale Liftung/Übereinstimmung/Aufgabe|Aufgabe 25.21]].

\inputaufgabe
{}
{

Es sei
\maabb {p} { E } { X
} {}
ein
\definitionsverweis {Vektorbündel}{}{}
vom
\definitionsverweis {Rang}{}{}
$r$ über einer
\definitionsverweis {differenzierbaren Mannigfaltigkeit}{}{}
$X$, das mit einem
\definitionsverweis {global integrablen Zusammenhang}{}{}
versehen sei. Zeige, dass es
\zusatzklammer {auf $X$ definierte} {} {}
\definitionsverweis {horizontale Schnitte}{}{}

\mathl{s_1  , \ldots , s_r}{} in $E$ derart gibt, dass
\maabbeledisp {} {X \times \R^r} { E
} {(P,v_1  , \ldots , v_r)} { \sum_{i = 1}^r v_i s_i(P)
} {,}
ein
\definitionsverweis {Isomorphismus von Vektorbündeln}{}{}
ist.

}
{} {}

\inputaufgabegibtloesung
{}
{

Es sei
\maabb {p} { E } { X
} {}
ein
\definitionsverweis {Vektorbündel}{}{}
vom
\definitionsverweis {Rang}{}{}
$r$ über einer
\definitionsverweis {differenzierbaren Mannigfaltigkeit}{}{}
$X$, das mit einem linearen
\definitionsverweis {global integrablen Zusammenhang}{}{}
versehen sei. Zeige, dass es
\zusatzklammer {auf $X$ definierte} {} {}
\definitionsverweis {horizontale Schnitte}{}{}

\mathl{s_1 , \ldots , s_r}{} in $E$ derart gibt, dass unter dem
\definitionsverweis {Isomorphismus von Vektorbündeln}{}{}
\maabbeledisp {} {X \times \R^r} { E
} {(P,v_1 , \ldots , v_r)} { \sum_{i = 1}^r v_i s_i(P)
} {,}
die
\definitionsverweis {Christoffelsymbole}{}{}
für den Zusammenhang bezüglich der Basisschnitte
\mathl{s_1  , \ldots , s_r}{} trivial sind.

}
{} {}

\inputaufgabegibtloesung
{}
{

Es sei ein zeitabhängiges stetiges
\definitionsverweis {Vektorfeld}{}{}
\maabbeledisp {F} {I \times \R^n } { \R^n
} { \left(  t   , \,  u_1  , \,   \ldots , \,  u_n                \right) } { \left(  F_1 \left(  u_1 , \ldots ,  u_n  \right)    , \,   \ldots  , \,   F_n \left(  u_1 , \ldots , u_n  \right)                 \right)
} {,}
mit dem zugehörigen
\definitionsverweis {gewöhnlichen Differentialgleichungssystem}{}{}

\mavergleichskettedisp
{\vergleichskette
{ u'
}
{ =} { F(t,u)
}
{ } {
}
{ } {
}
{ } {
}
}
{}{}{}
auf einem
\definitionsverweis {offenen Intervall}{}{}

\mavergleichskette
{\vergleichskette
{ I
}
{ \subseteq }{ \R
}
{  }{
}
{    }{
}
{   }{
}
}
{}{}{}
gegeben.
\aufzaehlungdrei{Zeige, dass durch
\maabbeledisp {} {I \times \R^n \times \R \times \R^n } { I \times \R^n  \times \R^n
} {(t,u,v,w)} { \left(  t   , \,  u  , \,   - vF \left(  t, u_1  , \ldots , u_n  \right)  +w                 \right)
} {,}
ein
\zusatzklammer {die vertikale Projektion für einen} {} {}
\definitionsverweis {Zusammenhang}{}{}
auf dem trivialen Vektorbündel
\maabbdisp {} {I \times \R^n } { \R^n
} {}
gegeben ist.
}{Bestimme die
\definitionsverweis {vertikale Ableitung}{}{}
$\nabla_\partial u$ zu einer
\definitionsverweis {differenzierbaren Kurve}{}{}
\maabbdisp {u} { I } { \R^n
} {}
\zusatzklammer {aufgefasst als Schnitt in $I \times \R^n$} {} {.}
}{Zeige, dass eine differenzierbare Kurve
\maabb {u} { I } { \R^n
} {}
genau dann eine Lösung des Differentialgleichungssystems ist, wenn $u$ ein
\definitionsverweis {horizontaler Schnitt}{}{}
bezüglich des Zusammenhangs ist.
}

}
{} {}

\inputaufgabe
{}
{

Es sei  $\nabla$ ein
\definitionsverweis {linearer Zusammenhang}{}{}
auf einem
\definitionsverweis {differenzierbaren Vektorbündel}{}{}
$E$ über einer
\definitionsverweis {differenzierbaren Mannigfaltigkeit}{}{}
$X$. Zeige, dass der Nullschnitt
\definitionsverweis {horizontal}{}{}
ist.

}
{} {}

\inputaufgabe
{}
{

Es sei $\nabla$ der
\definitionsverweis {triviale Zusammenhang}{}{}
auf dem trivialen Vektorbündel
\maabb {p} {X \times W } { X
} {}
zu einem reellen Vektorraum
\mathl{W}{} über einer
\definitionsverweis {differenzierbaren Mannigfaltigkeit}{}{}
$X$. Zeige, dass $\nabla$
\definitionsverweis {linear}{}{}
ist.

}
{} {}

\inputaufgabegibtloesung
{}
{

Es sei

\mavergleichskette
{\vergleichskette
{ U
}
{ \subseteq }{ \R^n
}
{  }{
}
{    }{
}
{   }{
}
}
{}{}{}
eine
\definitionsverweis {offene Teilmenge}{}{}
und es sei $\nabla$ der
\definitionsverweis {triviale Zusammenhang}{}{}
auf
\mathl{U \times \R}{.} Zeige, dass die
\definitionsverweis {Christoffelsymbole}{}{}
und die
\definitionsverweis {vertikale Ableitung}{}{}
des Zusammenhangs folgende Eigenschaften erfüllen.
\aufzaehlungvier{Die Christoffelsymbole sind

\mavergleichskettedisp
{\vergleichskette
{ \Gamma_{ij}^k
}
{ =} { 0
}
{ } {
}
{ } {
}
{ } {
}
}
{}{}{}
für alle $i,j,k$.
}{Es ist

\mavergleichskettedisp
{\vergleichskette
{ \nabla_{\partial_i }  \partial_j
}
{ =} { 0
}
{ } {
}
{ } {
}
{ } {
}
}
{}{}{}
für die
\definitionsverweis {Standardvektorfelder}{}{}
$\partial_i$.
}{Es ist

\mavergleichskettedisp
{\vergleichskette
{ \nabla_{\partial_i }  \left(  \sum_{j = 1}^n f_j \partial_j  \right)
}
{ =} { \sum_{j = 1}^n \partial_i(f_j) \partial_j
}
{ } {
}
{ } {
}
{ } {
}
}
{}{}{}
für differenzierbare Funktionen $f_j$.
}{Zu einem Vektorfeld $V$ und einem differenzierbaren Vektorfeld $W$ auf $U$ ist

\mavergleichskettedisp
{\vergleichskette
{ ( \nabla_V W)(P)
}
{ =} { \left(  D W   \right)_{ P } \left(   V(P)  \right)
}
{ } {
}
{ } {
}
{ } {
}
}
{}{}{.}
}

}
{} {}

\inputaufgabegibtloesung
{}
{

Wir betrachten auf dem
\definitionsverweis {trivialen Vektorbündel}{}{}

\mathl{\R \times \R^2}{} über $\R$ den
\definitionsverweis {trivialen Zusammenhang}{}{}
$\nabla$.
\aufzaehlungdrei{Zeige, dass

\mavergleichskette
{\vergleichskette
{ f_1(t)
}
{ = }{ \left(  1+ t^2  , \,  t^3                   \right)
}
{  }{
}
{    }{
}
{   }{
}
}
{}{}{}
und

\mavergleichskette
{\vergleichskette
{ f_2(t)
}
{ = }{ \left(  t  , \,   1+t^2                   \right)
}
{  }{
}
{    }{
}
{   }{
}
}
{}{}{}
\definitionsverweis {Basisschnitte}{}{}
des Vektorbündels sind.
}{Drücke die Standardbasisschnitte $e_1,e_2$ als Linearkombination der Basisschnitte $f_1,f_2$ aus.
}{Bestimme die
\definitionsverweis {Christoffelsymbole}{}{}
von $\nabla$ bezüglich dieser Basisschnitte.
}

}
{} {}

\inputaufgabe
{}
{

Wir betrachten auf dem
\definitionsverweis {trivialen Vektorbündel}{}{}

\mathl{\R^2 \times \R^2}{} über $\R^2$ den
\definitionsverweis {trivialen Zusammenhang}{}{}
$\nabla$.
\aufzaehlungzwei {Zeige, dass

\mavergleichskette
{\vergleichskette
{ f_1(x,y)
}
{ = }{ \left(  x   , \,   y                   \right)
}
{  }{
}
{    }{
}
{   }{
}
}
{}{}{}
und

\mavergleichskette
{\vergleichskette
{ f_2(x,y)
}
{ = }{ \left(  - y  , \,   x                   \right)
}
{  }{
}
{    }{
}
{   }{
}
}
{}{}{}
\definitionsverweis {Basisschnitte}{}{}
des Vektorbündels auf

\mavergleichskette
{\vergleichskette
{ U
}
{ = }{ \R^2 \setminus \{ (0,0)\}
}
{  }{
}
{    }{
}
{   }{
}
}
{}{}{} sind.
} {Bestimme die
\definitionsverweis {Christoffelsymbole}{}{}
von $\nabla$ bezüglich dieser Basisschnitte auf $U$.
}

}
{} {}

\inputaufgabegibtloesung
{}
{

Es sei ein
\definitionsverweis {homogenes lineares Differentialgleichungssystem}{}{}

\mavergleichskettedisp
{\vergleichskette
{ u'
}
{ =} { Mu
}
{ } {
}
{ } {
}
{ } {
}
}
{}{}{}
auf einem
\definitionsverweis {offenen Intervall}{}{}

\mavergleichskette
{\vergleichskette
{ I
}
{ \subseteq }{ \R
}
{  }{
}
{    }{
}
{   }{
}
}
{}{}{}
mit

\mavergleichskettedisp
{\vergleichskette
{ M(t)
}
{ =} { \left(  a_{kj} (t)  \right)_{kj}
}
{ } {
}
{ } {
}
{ } {
}
}
{}{}{}
\zusatzklammer {dabei ist $k$ der Zeilenindex und $j$ der Spaltenindex} {} {}
mit stetigen Funktionen
\maabbdisp {a_{kj}} { I } {\R
} {}
gegeben. Diese Funktionen fassen wir über

\mavergleichskettedisp
{\vergleichskette
{ \Gamma_{j}^k (t)
}
{ \defeq} { -a_{kj}
}
{ } {
}
{ } {
}
{ } {
}
}
{}{}{}
als
\definitionsverweis {Christoffelsymbole}{}{}
\zusatzklammer {mit

\mavergleichskettek
{\vergleichskettek
{ i
}
{ = }{ 1
}
{  }{
}
{    }{
}
{   }{
}
}
{}{}{}} {} {}
für den
\definitionsverweis {linearen Zusammenhang}{}{}
$\nabla$ auf dem trivialen Vektorbündel
\maabb {p} {I \times \R^n } { I
} {}
im Sinne von
[[Offene Menge/Triviales Vektorbündel/Linearer Zusammenhang/Spaltung/Christoffelsymbole/Bemerkung|Bemerkung 25.8]]
auf. Zeige, dass eine
\definitionsverweis {differenzierbare Kurve}{}{}
\maabbdisp {u} { I } { \R^n
} {}
genau dann eine Lösung des Differentialgleichungssystems ist, wenn $u$ ein
\definitionsverweis {horizontaler Schnitt}{}{}
bezüglich des Zusammenhangs ist.

}
{} {}

\inputaufgabe
{}
{

Es sei

\mavergleichskette
{\vergleichskette
{ Y
}
{ \subseteq }{ W
}
{ \subseteq }{ \R^n
}
{    }{
}
{   }{
}
}
{}{}{,}
$W$
\definitionsverweis {offen}{}{,}
eine
\definitionsverweis {differenzierbare Hyperfläche}{}{}
und sei
\maabb {\gamma} { I } { Y
} {}
eine
\definitionsverweis {differenzierbare Kurve}{}{.}
Es sei
\maabbdisp {F} { I } { TY
} {}
ein differenzierbares Vektorfeld längs $\gamma$, d.h. es liegt ein kommutatives Diagramm

\mathdisp {\begin{matrix}  &   &   & TY  \\ & & F \nearrow & \downarrow  \\  &  I  & \stackrel{ \gamma }{\longrightarrow}  &  Y \! \\ \end{matrix}} {  }

vor. Zeige, dass $F$ genau dann
\definitionsverweis {parallel längs}{}{}
$\gamma$ ist, wenn $F$ ein
\definitionsverweis {horizontaler Schnitt}{}{}
bezüglich des
\zusatzklammer {mithilfe der orthogonalen Projektion definierten} {} {}
\definitionsverweis {Zusammenhangs}{}{}
auf $TY$ ist.

}
{} {}

\inputaufgabe
{}
{

Es sei

\mavergleichskette
{\vergleichskette
{ Y
}
{ \subseteq }{ W
}
{ \subseteq }{ \R^n
}
{    }{
}
{   }{
}
}
{}{}{,}
$W$
\definitionsverweis {offen}{}{,}
eine
\definitionsverweis {differenzierbare Hyperfläche}{}{}
und sei $\nabla$ der
\zusatzklammer {mithilfe der orthogonalen Projektion definierte} {} {}
\definitionsverweis {Zusammenhang}{}{}
auf $TY$. Zeige, dass $\nabla$
\definitionsverweis {linear}{}{}
ist.

}
{} {}

\inputaufgabe
{}
{

Es seien
\mathkor {} {\nabla_1} {und} {\nabla_2} {}
\definitionsverweis {lineare Zusammenhänge}{}{}
auf einem
\definitionsverweis {differenzierbaren Vektorbündel}{}{}
$E$ über einer
\definitionsverweis {differenzierbaren Mannigfaltigkeit}{}{}
$X$. Es seien
\maabbdisp {g_1,g_2} { X } { \R
} {}
\definitionsverweis {stetig differenzierbare}{}{}
Funktionen mit

\mavergleichskette
{\vergleichskette
{ g_1+g_2
}
{ = }{ 1
}
{  }{
}
{    }{
}
{   }{
}
}
{}{}{.}
Zeige, dass
\mathl{g_1 \nabla_1+ g_2 \nabla_2}{} ebenfalls ein
\zusatzklammer {wie definierter} {?} {}
linearer  Zusammenhang auf $E$  ist.

}
{} {}

\inputaufgabe
{}
{

Es sei

\mavergleichskette
{\vergleichskette
{ X
}
{ = }{ \bigcup_{i \in I} U_i
}
{  }{
}
{    }{
}
{   }{
}
}
{}{}{}
eine
\definitionsverweis {offene Überdeckung}{}{}
einer
\definitionsverweis {differenzierbaren Mannigfaltigkeit}{}{}
$X$, auf der das
\definitionsverweis {differenzierbare Vektorbündel}{}{}
$E$ trivialisiert. Es sei  $\nabla_i$ der
\definitionsverweis {triviale Zusammenhang}{}{}
auf $E\vert_{U_i }$ und es sei
\mathkor {} {h_j} {} {j \in J} {,}
eine der Überdeckung untergeordnete
\definitionsverweis {Partition der Eins}{}{,}

\mathbed {h_j} {}
 {j \in J} {}
 {} {} {} {.}
Zeige, dass $\sum_{j \in J} h_j \nabla_{i(j)}$ ein wohldefinierter linearer Zusammenhang auf $E$ ist.

}
{} {}

\inputaufgabe
{}
{

Es sei $X$ eine
\definitionsverweis {differenzierbare Mannigfaltigkeit}{}{}
und sei $\omega$ eine stetige
$1$-\definitionsverweis {Differentialform}{}{}
auf $X$. Zeige, dass die Abbildung
\maabbeledisp {} {C^1(X,\R)} { C^0(X,T^*X)
} { f } { \nabla_\omega (f) \defeq f \omega + df
} {,}
die Leibniz-Regel aus
[[Mannigfaltigkeit/Vektorbündel/R/Linearer Zusammenhang/Leibnizregel/Fakt|Satz 25.5]]
erfüllt.

}
{} {}

\inputaufgabe
{}
{

Es sei
\maabb {p} { E } { X
} {}
ein differenzierbares
\definitionsverweis {Vektorbündel}{}{}
über einer
\definitionsverweis {differenzierbaren Mannigfaltigkeit}{}{}
$X$. Es sei $E$ selbst eine
\definitionsverweis {riemannsche Mannigfaltigkeit}{}{.}
Zeige, dass auf $E$ ein
\definitionsverweis {linearer Zusammenhang}{}{}
gegeben ist, indem man zum
\definitionsverweis {Vertikalbündel}{}{}

\mavergleichskette
{\vergleichskette
{ V
}
{ \subseteq }{ T E
}
{  }{
}
{    }{
}
{   }{
}
}
{}{}{}
über das
\definitionsverweis {orthogonale Komplement}{}{}
ein
\definitionsverweis {Horizontalbündel}{}{}
definiert.

}
{} {}

\inputaufgabe
{}
{

Es sei
\maabb {p} { E } { X
} {}
ein
\definitionsverweis {Vektorbündel}{}{}
über der
\definitionsverweis {differenzierbaren Mannigfaltigkeit}{}{}
$X$ mit
\definitionsverweis {abzählbarer Basis der Topologie}{}{.}
Zeige mit
[[Mannigfaltigkeit/Abzählbare Basis/Vektorbündel/Riemannsche Struktur/Aufgabe|Aufgabe 22.15]]
und mit
[[Mannigfaltigkeit/Vektorbündel/Riemannsche Mannigfaltigkeit/Linearer Zusammenhang/Aufgabe|Aufgabe 25.18]],
dass es auf $E$ einen
\definitionsverweis {linearen Zusammenhang}{}{}
gibt.

}
{} {}

  \zwischenueberschrift{Aufgaben zum Abgeben}

\inputaufgabe
{5 (3+2)}
{

Es sei

\mavergleichskette
{\vergleichskette
{ X
}
{ = }{ S^1
}
{  }{
}
{    }{
}
{   }{
}
}
{}{}{}
der
\definitionsverweis {Einheitskreis}{}{}
mit dem
\definitionsverweis {Tangentialbündel}{}{}
\maabb {p} {TS^1 = S^1 \times \R} { S^1
} {.}
\aufzaehlungzwei {Beschreibe einen
\definitionsverweis {Zusammenhang}{}{}
auf $TS^1$ derart, dass es keinen
\definitionsverweis {horizontalen Schnitt}{}{}
auf ganz $S^1$ gibt.
} {Bestimme einen horizontalen Schnitt zum Zusammenhang aus (1) längs der Abbildung
\maabbeledisp {} {\R} { S^1
} { t } { \left(  \cos  t   , \,   \sin  t                   \right)
} {.}
}

}
{} {}

\inputaufgabe
{4}
{

Es sei $X$ eine
\definitionsverweis {differenzierbare Mannigfaltigkeit}{}{}
und $E$ ein
\definitionsverweis {differenzierbares Vektorbündel}{}{}
auf $X$, das mit einem
\definitionsverweis {Zusammenhang}{}{}
versehen sei. Es sei $Z$ eine weitere differenzierbare Mannigfaltigkeit mit einer differenzierbaren Abbildung
\maabb {\psi} { Z } { X
} {.}
Es sei $Z$
\definitionsverweis {zusammenhängend}{}{}
und seien
\maabb {s_1,s_2} { Z } { E
} {}
\definitionsverweis {horizontale Schnitte}{}{}
über $\psi$ mit

\mavergleichskette
{\vergleichskette
{ s_1(z)
}
{ = }{ s_2(z)
}
{  }{
}
{    }{
}
{   }{
}
}
{}{}{}
für einen Punkt

\mavergleichskette
{\vergleichskette
{ z
}
{ \in }{ Z
}
{  }{
}
{    }{
}
{   }{
}
}
{}{}{.}
Zeige

\mavergleichskette
{\vergleichskette
{ s_1
}
{ = }{ s_2
}
{  }{
}
{    }{
}
{   }{
}
}
{}{}{.}

}
{} {}

\inputaufgabe
{3}
{

Es sei $X$ eine
\definitionsverweis {differenzierbare Mannigfaltigkeit}{}{}
und $E$ ein
\definitionsverweis {differenzierbares Vektorbündel}{}{}  auf $X$, das mit einem
\definitionsverweis {linearen Zusammenhang}{}{}
versehen sei und es sei $G$ ein
\definitionsverweis {Vektorfeld}{}{}
auf $X$. Zeige, dass für jeden
\definitionsverweis {differenzierbaren Schnitt}{}{}
$s$ in $E$ und jede
\definitionsverweis {differenzierbare Funktion}{}{}
$f$
\zusatzklammer {die beide auf einer offenen Menge

\mavergleichskette
{\vergleichskette
{ U
}
{ \subseteq }{ X
}
{  }{
}
{    }{
}
{   }{
}
}
{}{}{}
definiert sind} {} {}
die Beziehung

\mavergleichskettedisp
{\vergleichskette
{ \nabla_G (fs)
}
{ =} { s \otimes D_G (f) + f \nabla_G (s)
}
{ } {
}
{ } {
}
{ } {
}
}
{}{}{}
gilt.

}
{} {}

\inputaufgabe
{3}
{

Man gebe ein Beispiel für einen
\definitionsverweis {linearen Zusammenhang}{}{}
auf einem
\definitionsverweis {Vektorbündel}{}{}
\maabb {p} { E } { X
} {}
über einer
\definitionsverweis {differenzierbaren Mannigfaltigkeit}{}{}
$X$ derart, dass der Vektorraum
\mathl{H(\nabla,X)}{}  der
\definitionsverweis {horizontalen Schnitte}{}{}
eine
\definitionsverweis {Vektorraumdimension}{}{}
besitzt, die größer als der
\definitionsverweis {Rang}{}{}
von $E$ ist.

}
{} {}

 [[Kategorie:Latexseite]]
